\section*{Task 0: Projectives and syzygies}

\paragraph{Goal.}
Compute projective covers, syzygies, and the basic \(\Ext^1\)-criterion for the
weight-space cyclic Nakayama algebra.

\paragraph{Projective indecomposables in lifted coordinates.}
In the lifted indexing of Conjecture 3, the indecomposable projectives are
\[
P(0)=M(0,L),
\qquad
P(a)=M(a,a+L-1)\quad (1\le a\le L).
\]
Equivalently, the maximal admissible endpoint for the start vertex \(a\) is
\[
\beta(a):=
\begin{cases}
L, & a=0,\\
a+L-1, & 1\le a\le L.
\end{cases}
\]

\paragraph{Projective cover and first syzygy.}
For any nonprojective indecomposable \(M(a,b)\), the projective cover is
\[
P(a)\twoheadrightarrow M(a,b),
\]
and its kernel is the interval module
\[
\Omega M(a,b)\cong M(b+1,\beta(a)).
\]
Thus there is a short exact sequence
\[
0\longrightarrow M(b+1,\beta(a))
\longrightarrow P(a)
\longrightarrow M(a,b)
\longrightarrow 0.
\]
This is immediate from the uniserial structure of the projective \(P(a)\): the
quotient by the first \(b-a+1\) composition factors is exactly \(M(a,b)\), so
the remaining tail is \(M(b+1,\beta(a))\).

\paragraph{Hom criterion for interval modules.}
For admissible lifted intervals \((a,b)\) and \((c,d)\),
\[
\dim_k \Hom_A(M(a,b),M(c,d))
=
\begin{cases}
1, & [a,b]\subseteq [c,d],\\
0, & \text{otherwise}.
\end{cases}
\]
Indeed, every indecomposable interval module has pairwise distinct composition
factors, so any nonzero morphism is determined up to scalar and exists exactly
when the source interval embeds into the target interval.

\paragraph{Ext criterion.}
Combining the previous two facts yields
\[
\Ext_A^1(M(a,b),M(c,d))
\cong
\Hom_A(M(b+1,\beta(a)),M(c,d))
\]
for every nonprojective \(M(a,b)\). Therefore
\[
\dim_k \Ext_A^1(M(a,b),M(c,d))
=
\begin{cases}
1, & [b+1,\beta(a)]\subseteq [c,d],\\
0, & \text{otherwise},
\end{cases}
\]
and the dimension is \(0\) when \(M(a,b)\) is projective.

\paragraph{Status.}
Completed. The homological structure is explicit and can be used either for
\(\Ext^1\)-based deformation arguments or for comparison with the older draft.
